\section{Classical BCA Prototype Model}

\subsection{Preferences and Technology}

The representative household has expected utility
\begin{equation}
    \E_0
    \sum_{t=0}^{\infty}
    \beta^t U(C_t,L_t),
    \qquad 0<\beta<1.
    \label{eq:utility}
\end{equation}
For algebraic examples, one can use
\begin{equation}
    U(C_t,L_t)
    =
    \frac{C_t^{1-\gamma}-1}{1-\gamma}
    -
    \chi\frac{L_t^{1+\nu}}{1+\nu},
    \qquad
    \gamma>0,\ \nu>0,\ \chi>0.
    \label{eq:period_utility_example}
\end{equation}

Output is produced with capital, labor, and total factor productivity:
\begin{equation}
    Y_t
    =
    A_t F(K_t,L_t)
    =
    A_t K_t^\alpha L_t^{1-\alpha},
    \qquad 0<\alpha<1.
    \label{eq:production}
\end{equation}
This is the location of the efficiency wedge. In the prototype model,
\(A_t\) is not merely a Solow residual; it is the object that absorbs any
distortion that makes measured output differ from what capital and labor alone
would predict.

\subsection{Resource Constraint and Capital Accumulation}

The closed-economy resource constraint is
\begin{equation}
    C_t + X_t + G_t = Y_t,
    \label{eq:resource}
\end{equation}
where \(X_t\) is investment and \(G_t\) is the government or resource wedge. In
BCA, \(G_t\) should be interpreted broadly: it includes government purchases,
net exports in open-economy variants, measurement residuals, and any other
resource absorption not modeled as consumption or investment.

Capital evolves according to
\begin{equation}
    K_{t+1}
    =
    (1-\delta)K_t + X_t,
    \qquad 0<\delta<1.
    \label{eq:capital_accumulation}
\end{equation}

\subsection{Canonical Wedges}

The labor wedge is defined by the intratemporal condition
\begin{equation}
    -\frac{U_L(C_t,L_t)}{U_C(C_t,L_t)}
    =
    (1-\tau_{\ell,t})A_t F_L(K_t,L_t).
    \label{eq:labor_wedge}
\end{equation}
When \(\tau_{\ell,t}=0\), the household's marginal rate of substitution equals
the marginal product of labor. A positive labor wedge lowers the effective
return to working.

Define the gross marginal product return on capital as
\begin{equation}
    R^k_{t+1}
    =
    A_{t+1}F_K(K_{t+1},L_{t+1}) + 1-\delta.
    \label{eq:capital_return}
\end{equation}
The investment wedge is defined by
\begin{equation}
    1
    =
    \beta
    \E_t
    \left[
        \frac{U_C(C_{t+1},L_{t+1})}{U_C(C_t,L_t)}
        (1-\tau_{x,t})
        R^k_{t+1}
    \right].
    \label{eq:investment_wedge}
\end{equation}
This normalization places the investment wedge at date \(t\) as a tax on the
expected return to saving through capital. Other BCA normalizations place the
wedge as a tax on investment goods or index it at \(t+1\). These conventions are
equivalent up to timing and sign conventions as long as the definition is held
fixed.

The four canonical wedges are therefore:
\begin{align}
    \text{Efficiency wedge:} \quad & A_t,
    \label{eq:efficiency_wedge_def}
    \\
    \text{Government/resource wedge:} \quad & G_t = Y_t-C_t-X_t,
    \label{eq:government_wedge_def}
    \\
    \text{Labor wedge:} \quad &
    1-\tau_{\ell,t}
    =
    \frac{-U_L(C_t,L_t)}
    {U_C(C_t,L_t)A_tF_L(K_t,L_t)},
    \label{eq:labor_wedge_def}
    \\
    \text{Investment wedge:} \quad &
    1-\tau_{x,t}
    =
    \left\{
    \beta
    \E_t
    \left[
        \frac{U_C(C_{t+1},L_{t+1})}{U_C(C_t,L_t)}
        R^k_{t+1}
    \right]
    \right\}^{-1}.
    \label{eq:investment_wedge_def}
\end{align}

\begin{remark}
There is no Phillips curve, markup wedge, or Taylor-rule wedge in this
prototype. If one later wants to study price rigidity, it should be introduced
as a structural mechanism that may show up through the labor wedge or other
canonical wedges, not as an additional BCA wedge.
\end{remark}
